% =====================================================================
% Introduction skeleton for the ISCA 2027 µPyCore paper.
%
% Drop-in:   % =====================================================================
% Introduction skeleton for the ISCA 2027 µPyCore paper.
%
% Drop-in:   % =====================================================================
% Introduction skeleton for the ISCA 2027 µPyCore paper.
%
% Drop-in:   % =====================================================================
% Introduction skeleton for the ISCA 2027 µPyCore paper.
%
% Drop-in:   \input{sections/introduction}  (replaces the placeholder
%            \section{Introduction} in the main file)
% Citations: keys resolve against docs/paper/references.bib. To use it,
%            replace the template thebibliography block with
%              \bibliographystyle{IEEEtran}
%              \bibliography{references}
%
% Conventions used below
%   % P<n> GOAL:   what the paragraph must accomplish for a reviewer
%   % EVIDENCE:    which citations carry the claim
%   \textcolor{red}{...}   number still to be finalized (matches the
%                          convention in the Evaluation sections)
%   % TODO:        prose or a fact still to settle
%
% Target length: ~1 to 1.25 columns of text + the contribution list.
% ISCA reviewers read the intro to answer four questions, in order:
%   (1) why does this problem matter now, (2) why do existing
%   solutions fail, (3) what is the key insight, (4) what did you
%   show. Each paragraph below answers one of them.
% =====================================================================

\section{Introduction}
\label{sec:intro}

% ---------------------------------------------------------------------
% P1 GOAL: The "why now". AI coding has moved the bottleneck from
% writing code to reading and reviewing it. Keep it to hard numbers;
% reviewers at ISCA are skeptical of AI hype, so lead with measured
% effects rather than CEO quotes.
% EVIDENCE: pichai2024, nadella2025 (scale of AI-written code);
% faros2025 (review time); becker2025metr (RCT); stackoverflow2025
% (developer distrust -> more review).
% ---------------------------------------------------------------------
Large language models now write a substantial fraction of production
code: Google and Microsoft report that 25\% and 20--30\% of their new
code, respectively, is generated by AI~\cite{pichai2024,nadella2025}.
Writing code is no longer the bottleneck; checking it is. Across
roughly 10,000 developers, teams with high AI adoption merge 98\% more
pull requests, but each one spends 91\% longer in review~\cite{faros2025}.
In a randomized controlled trial, experienced open-source developers
took 19\% \emph{longer} to finish tasks with AI assistance, even though
they believed it had made them faster~\cite{becker2025metr}. Developers
also trust what they review less: 46\% now distrust the accuracy of AI
output, up from 31\% a year earlier~\cite{stackoverflow2025}.
% TODO: one sentence of the authors' thesis: when humans mainly review,
% the language should be optimized for the reader. Python was designed
% around exactly that principle ("Readability counts"~\cite{pep20}).

% ---------------------------------------------------------------------
% P2 GOAL: Python is where human preference and model capability
% converge, so it is the language worth building hardware for. This
% paragraph justifies the choice of Python specifically (a reviewer
% will ask "why not JavaScript / Lua / Java?").
% EVIDENCE: octoverse2024, stackoverflow2025 (humans); twist2025,
% kocetkov2023, cassano2023 (models).
% ---------------------------------------------------------------------
Python is where both trends meet. In 2024 it passed JavaScript as the
most used language on GitHub~\cite{octoverse2024}, and its use among
surveyed developers grew by 7 percentage points in a single
year~\cite{stackoverflow2025}. LLMs show the same preference. Asked to
solve benchmark tasks without being told a language, eight models
chose Python in 90--97\% of cases, and still chose it in 58\% of project
setups where it was a poor fit~\cite{twist2025}. Python is
also heavily represented in the open code corpora these models are
trained on~\cite{kocetkov2023}, and code-generation quality varies
with how well a language is represented in that data~\cite{cassano2023}.
% TODO(verify): before submission, check the direction and size of the
% MultiPL-E popularity/accuracy finding in the paper itself. Search
% snippets confirm only that the paper studies it.

% ---------------------------------------------------------------------
% P3 GOAL: The cost. Python is slow, and the gap is worst exactly
% where a JIT cannot help: memory-constrained embedded devices. This is
% the problem statement; end it with a crisp sentence that a reviewer
% can quote back.
% EVIDENCE: leiserson2020 (47x language gap), hennessy2019 (frames the
% Python gap as an architecture opportunity), ismail2018 (where the
% interpreter's time goes), pereira2021 (energy), micropython,
% eclipse2024 (Python on constrained devices).
% ---------------------------------------------------------------------
That readability comes at a steep cost. Rewriting a Python matrix
multiply in C makes it 47$\times$ faster before any hardware-specific
tuning~\cite{leiserson2020}. In energy-efficiency rankings across 27
languages, Python places near the bottom~\cite{pereira2021}. Hennessy and
Patterson identify this gap as one of the main opportunities for
computer architects~\cite{hennessy2019}. On servers and laptops, JIT
compilers recover part of the gap by spending memory and warm-up
time~\cite{bolz2009pypy,pep744}. Microcontrollers cannot afford either.
MicroPython~\cite{micropython} brings Python to devices with hundreds of
kilobytes of RAM, and 32\% of developers targeting constrained devices
already use Python~\cite{eclipse2024}. Yet these devices run Python
through a bytecode interpreter, where dynamic type dispatch, object
allocation and calls into the C runtime dominate execution
time~\cite{ismail2018}.
% TODO: optionally add one sentence on the size of the MCU market
% (see the yole2025 note in references.bib; a sourced unit count is
% still missing).
% TODO: close P3 with a one-line problem statement, e.g. "Embedded
% developers must choose between the language that people and models
% read best and the performance their device needs."

% ---------------------------------------------------------------------
% P4 GOAL: Why prior work does not solve it. Three buckets, one or
% two sentences each. Each sentence should end with what that bucket
% gives up, which sets up the insight in P5. The detailed versions live
% in Sec. Background / Related Work, so point there.
% EVIDENCE: pep659, pep744, bolz2009pypy (software); micropython
% (native/viper emitters, see Background); mcghan1998, schoeberl2008,
% porthouse2005 (Java processors); moon1985 (tagged Lisp machine);
% kim2017typed, choi2017shortcut, gope2017php, maas2018gc (assists for
% conventional cores); pyxl2025.
% ---------------------------------------------------------------------
Existing approaches recover performance by giving something up.
\emph{Software} techniques such as CPython's specializing
interpreter~\cite{pep659} and tracing or copy-and-patch
JITs~\cite{bolz2009pypy,pep744} need memory a microcontroller lacks.
MicroPython's own native and Viper emitters instead give up the managed
semantics that make the language worth using (Section~\ref{sec:background}).
% TODO: \label{sec:background} does not exist yet; add it to the
% Background section header.
\emph{Language processors} executed bytecode directly, but they targeted
statically typed, verifiable Java bytecode~\cite{mcghan1998,schoeberl2008,porthouse2005}.
They kept the interpreter's hard cases, such as object operations, in
software. Earlier tagged Lisp machines~\cite{moon1985} showed that
type tags in hardware help a dynamic language. \emph{Architectural
assists} for scripting languages~\cite{kim2017typed,choi2017shortcut,gope2017php,maas2018gc}
add type checks, dispatch shortcuts, or hash and GC
engines to a conventional core that still runs an interpreter or JIT.
The reported gains are 10--30\%.
% TODO(verify): the 10-30% range combines kim2017typed (about 10-11%
% geomean) and choi2017shortcut (30% under the baseline compiler).
% Check both against the papers, or drop the number.
Most recently, PyXL~\cite{pyxl2025} runs Python on an FPGA, but it
compiles to a custom instruction set on a host machine and is not
publicly documented.

% ---------------------------------------------------------------------
% P5 GOAL: The key insight + what µPyCore is. State the thesis in one
% sentence, then map each obstacle (from the Background "rules") to the
% mechanism that removes it. This is the paragraph reviewers summarize
% in their review, so make the mapping explicit.
% ---------------------------------------------------------------------
This paper takes the opposite approach: rather than lowering Python
toward the machine, we raise the machine to Python. \textmu{}PyCore
executes unmodified MicroPython bytecode as its native instruction set
and implements in hardware the runtime machinery that a conventional
core spends most of its time emulating. Three obstacles shape the
design.
% TODO: tighten each to one sentence once Sec. Design is final.
First, bytecode carries no type information, so every value in
\textmu{}PyCore carries a 4-bit type tag that the core reads together
with its payload in a single access (Section~\ref{sec:tags}).
Second, the bytecode targets a stack machine, so the register file is a
tagged ring buffer that holds the top of the execution stack and
spills to SRAM in fixed frames (Section~\ref{sec:regfile}).
Third, the language's time goes to built-in containers, allocation and
garbage collection rather than to arithmetic. A collection accelerator
and two range-based scoreboards therefore run that work in parallel with
the core (Sections~\ref{sec:accelerator}--\ref{sec:builtins}).
% TODO: Scoreboarding subsection needs a \label (e.g. sec:scoreboard).
Everything the hardware does not cover, including \texttt{print},
\texttt{sorted} and a full bytecode compiler, runs as MicroPython
firmware on the same core, so Python source can be compiled and run on
the device itself.

% ---------------------------------------------------------------------
% P6 GOAL: Headline results. Same numbers as the abstract; keep red
% until the evaluation is frozen. State the methodology caveat
% (RTL simulation at the post-P&R clock vs. published board numbers)
% in the same breath so a reviewer cannot say it was hidden.
% ---------------------------------------------------------------------
We synthesize and place-and-route \textmu{}PyCore in SkyWater's open
130-nm process~\cite{skywater130}, where it closes timing at
\textcolor{red}{83.33~MHz}. Cycle counts from RTL simulation at that
clock project \textcolor{red}{7,286}~pystones/s,
\textcolor{red}{2.84$\times$} the published throughput of an Arm-based
Raspberry Pi Pico~2 at 150~MHz and \textcolor{red}{3.19$\times$} that of a
PyBoard at 168~MHz. Across \textcolor{red}{15} workloads from
MicroPython's \texttt{perf\_bench} suite, \textmu{}PyCore achieves a
geometric-mean speedup of \textcolor{red}{2.33$\times$} over the fastest
published microcontroller result for each workload, while running at a
lower clock rate than any of them.
% TODO: reconcile the numbers with the abstract and Sec. Eval before
% submission (see the notes left with this skeleton).

% ---------------------------------------------------------------------
% Contributions. ISCA reviewers look for this list. Make each item
% something that can be checked, and point it at its section.
% ---------------------------------------------------------------------
This paper makes the following contributions:
\begin{itemize}
  \item A processor that runs MicroPython~v1.29.0 bytecode as its native ISA, covering
        \textcolor{red}{N} of the 207 defined opcodes with documented
        semantic deviations (Section~\ref{sec:bytecode-isa}).
  \item A type-tagged, ring-buffer register file that maps the
        interpreter's operand stack onto hardware and spills in fixed
        frames (Sections~\ref{sec:regfile}--\ref{sec:tags}).
  \item A collection accelerator and range-based memory and ALU scoreboards
        that overlap container upkeep and arithmetic with execution
        while keeping exceptions precise (Section~\ref{sec:accelerator}).
  \item A hardware/firmware split of the Python runtime, including builtins,
        standard-library modules, garbage collection and a self-hosted
        on-device compiler written in MicroPython (Section~\ref{sec:builtins}).
  \item A post-place-and-route Sky130 implementation evaluated on
        MicroPython's \texttt{perf\_bench} suite against published results
        for five commercial microcontroller configurations (Section~\textcolor{red}{X}).
  % TODO: decide whether the custom benchmarks earn their own bullet.
  % TODO: if the RTL will be open-sourced, say so here (reviewers value
  % artifacts); keep it anonymized for review.
\end{itemize}

% ---------------------------------------------------------------------
% Roadmap. Optional at ISCA; drop it if space is tight.
% ---------------------------------------------------------------------
% Section~\ref{sec:background} reviews how MicroPython executes bytecode
% and derives the requirements for running it natively.
% Section~\ref{sec:design} presents the design of \textmu{}PyCore.
% Sections~X and~Y describe our methodology and results, and
% Section~Z discusses related work.
  (replaces the placeholder
%            \section{Introduction} in the main file)
% Citations: keys resolve against docs/paper/references.bib. To use it,
%            replace the template thebibliography block with
%              \bibliographystyle{IEEEtran}
%              \bibliography{references}
%
% Conventions used below
%   % P<n> GOAL:   what the paragraph must accomplish for a reviewer
%   % EVIDENCE:    which citations carry the claim
%   \textcolor{red}{...}   number still to be finalized (matches the
%                          convention in the Evaluation sections)
%   % TODO:        prose or a fact still to settle
%
% Target length: ~1 to 1.25 columns of text + the contribution list.
% ISCA reviewers read the intro to answer four questions, in order:
%   (1) why does this problem matter now, (2) why do existing
%   solutions fail, (3) what is the key insight, (4) what did you
%   show. Each paragraph below answers one of them.
% =====================================================================

\section{Introduction}
\label{sec:intro}

% ---------------------------------------------------------------------
% P1 GOAL: The "why now". AI coding has moved the bottleneck from
% writing code to reading and reviewing it. Keep it to hard numbers;
% reviewers at ISCA are skeptical of AI hype, so lead with measured
% effects rather than CEO quotes.
% EVIDENCE: pichai2024, nadella2025 (scale of AI-written code);
% faros2025 (review time); becker2025metr (RCT); stackoverflow2025
% (developer distrust -> more review).
% ---------------------------------------------------------------------
Large language models now write a substantial fraction of production
code: Google and Microsoft report that 25\% and 20--30\% of their new
code, respectively, is generated by AI~\cite{pichai2024,nadella2025}.
Writing code is no longer the bottleneck; checking it is. Across
roughly 10,000 developers, teams with high AI adoption merge 98\% more
pull requests, but each one spends 91\% longer in review~\cite{faros2025}.
In a randomized controlled trial, experienced open-source developers
took 19\% \emph{longer} to finish tasks with AI assistance, even though
they believed it had made them faster~\cite{becker2025metr}. Developers
also trust what they review less: 46\% now distrust the accuracy of AI
output, up from 31\% a year earlier~\cite{stackoverflow2025}.
% TODO: one sentence of the authors' thesis: when humans mainly review,
% the language should be optimized for the reader. Python was designed
% around exactly that principle ("Readability counts"~\cite{pep20}).

% ---------------------------------------------------------------------
% P2 GOAL: Python is where human preference and model capability
% converge, so it is the language worth building hardware for. This
% paragraph justifies the choice of Python specifically (a reviewer
% will ask "why not JavaScript / Lua / Java?").
% EVIDENCE: octoverse2024, stackoverflow2025 (humans); twist2025,
% kocetkov2023, cassano2023 (models).
% ---------------------------------------------------------------------
Python is where both trends meet. In 2024 it passed JavaScript as the
most used language on GitHub~\cite{octoverse2024}, and its use among
surveyed developers grew by 7 percentage points in a single
year~\cite{stackoverflow2025}. LLMs show the same preference. Asked to
solve benchmark tasks without being told a language, eight models
chose Python in 90--97\% of cases, and still chose it in 58\% of project
setups where it was a poor fit~\cite{twist2025}. Python is
also heavily represented in the open code corpora these models are
trained on~\cite{kocetkov2023}, and code-generation quality varies
with how well a language is represented in that data~\cite{cassano2023}.
% TODO(verify): before submission, check the direction and size of the
% MultiPL-E popularity/accuracy finding in the paper itself. Search
% snippets confirm only that the paper studies it.

% ---------------------------------------------------------------------
% P3 GOAL: The cost. Python is slow, and the gap is worst exactly
% where a JIT cannot help: memory-constrained embedded devices. This is
% the problem statement; end it with a crisp sentence that a reviewer
% can quote back.
% EVIDENCE: leiserson2020 (47x language gap), hennessy2019 (frames the
% Python gap as an architecture opportunity), ismail2018 (where the
% interpreter's time goes), pereira2021 (energy), micropython,
% eclipse2024 (Python on constrained devices).
% ---------------------------------------------------------------------
That readability comes at a steep cost. Rewriting a Python matrix
multiply in C makes it 47$\times$ faster before any hardware-specific
tuning~\cite{leiserson2020}. In energy-efficiency rankings across 27
languages, Python places near the bottom~\cite{pereira2021}. Hennessy and
Patterson identify this gap as one of the main opportunities for
computer architects~\cite{hennessy2019}. On servers and laptops, JIT
compilers recover part of the gap by spending memory and warm-up
time~\cite{bolz2009pypy,pep744}. Microcontrollers cannot afford either.
MicroPython~\cite{micropython} brings Python to devices with hundreds of
kilobytes of RAM, and 32\% of developers targeting constrained devices
already use Python~\cite{eclipse2024}. Yet these devices run Python
through a bytecode interpreter, where dynamic type dispatch, object
allocation and calls into the C runtime dominate execution
time~\cite{ismail2018}.
% TODO: optionally add one sentence on the size of the MCU market
% (see the yole2025 note in references.bib; a sourced unit count is
% still missing).
% TODO: close P3 with a one-line problem statement, e.g. "Embedded
% developers must choose between the language that people and models
% read best and the performance their device needs."

% ---------------------------------------------------------------------
% P4 GOAL: Why prior work does not solve it. Three buckets, one or
% two sentences each. Each sentence should end with what that bucket
% gives up, which sets up the insight in P5. The detailed versions live
% in Sec. Background / Related Work, so point there.
% EVIDENCE: pep659, pep744, bolz2009pypy (software); micropython
% (native/viper emitters, see Background); mcghan1998, schoeberl2008,
% porthouse2005 (Java processors); moon1985 (tagged Lisp machine);
% kim2017typed, choi2017shortcut, gope2017php, maas2018gc (assists for
% conventional cores); pyxl2025.
% ---------------------------------------------------------------------
Existing approaches recover performance by giving something up.
\emph{Software} techniques such as CPython's specializing
interpreter~\cite{pep659} and tracing or copy-and-patch
JITs~\cite{bolz2009pypy,pep744} need memory a microcontroller lacks.
MicroPython's own native and Viper emitters instead give up the managed
semantics that make the language worth using (Section~\ref{sec:background}).
% TODO: \label{sec:background} does not exist yet; add it to the
% Background section header.
\emph{Language processors} executed bytecode directly, but they targeted
statically typed, verifiable Java bytecode~\cite{mcghan1998,schoeberl2008,porthouse2005}.
They kept the interpreter's hard cases, such as object operations, in
software. Earlier tagged Lisp machines~\cite{moon1985} showed that
type tags in hardware help a dynamic language. \emph{Architectural
assists} for scripting languages~\cite{kim2017typed,choi2017shortcut,gope2017php,maas2018gc}
add type checks, dispatch shortcuts, or hash and GC
engines to a conventional core that still runs an interpreter or JIT.
The reported gains are 10--30\%.
% TODO(verify): the 10-30% range combines kim2017typed (about 10-11%
% geomean) and choi2017shortcut (30% under the baseline compiler).
% Check both against the papers, or drop the number.
Most recently, PyXL~\cite{pyxl2025} runs Python on an FPGA, but it
compiles to a custom instruction set on a host machine and is not
publicly documented.

% ---------------------------------------------------------------------
% P5 GOAL: The key insight + what µPyCore is. State the thesis in one
% sentence, then map each obstacle (from the Background "rules") to the
% mechanism that removes it. This is the paragraph reviewers summarize
% in their review, so make the mapping explicit.
% ---------------------------------------------------------------------
This paper takes the opposite approach: rather than lowering Python
toward the machine, we raise the machine to Python. \textmu{}PyCore
executes unmodified MicroPython bytecode as its native instruction set
and implements in hardware the runtime machinery that a conventional
core spends most of its time emulating. Three obstacles shape the
design.
% TODO: tighten each to one sentence once Sec. Design is final.
First, bytecode carries no type information, so every value in
\textmu{}PyCore carries a 4-bit type tag that the core reads together
with its payload in a single access (Section~\ref{sec:tags}).
Second, the bytecode targets a stack machine, so the register file is a
tagged ring buffer that holds the top of the execution stack and
spills to SRAM in fixed frames (Section~\ref{sec:regfile}).
Third, the language's time goes to built-in containers, allocation and
garbage collection rather than to arithmetic. A collection accelerator
and two range-based scoreboards therefore run that work in parallel with
the core (Sections~\ref{sec:accelerator}--\ref{sec:builtins}).
% TODO: Scoreboarding subsection needs a \label (e.g. sec:scoreboard).
Everything the hardware does not cover, including \texttt{print},
\texttt{sorted} and a full bytecode compiler, runs as MicroPython
firmware on the same core, so Python source can be compiled and run on
the device itself.

% ---------------------------------------------------------------------
% P6 GOAL: Headline results. Same numbers as the abstract; keep red
% until the evaluation is frozen. State the methodology caveat
% (RTL simulation at the post-P&R clock vs. published board numbers)
% in the same breath so a reviewer cannot say it was hidden.
% ---------------------------------------------------------------------
We synthesize and place-and-route \textmu{}PyCore in SkyWater's open
130-nm process~\cite{skywater130}, where it closes timing at
\textcolor{red}{83.33~MHz}. Cycle counts from RTL simulation at that
clock project \textcolor{red}{7,286}~pystones/s,
\textcolor{red}{2.84$\times$} the published throughput of an Arm-based
Raspberry Pi Pico~2 at 150~MHz and \textcolor{red}{3.19$\times$} that of a
PyBoard at 168~MHz. Across \textcolor{red}{15} workloads from
MicroPython's \texttt{perf\_bench} suite, \textmu{}PyCore achieves a
geometric-mean speedup of \textcolor{red}{2.33$\times$} over the fastest
published microcontroller result for each workload, while running at a
lower clock rate than any of them.
% TODO: reconcile the numbers with the abstract and Sec. Eval before
% submission (see the notes left with this skeleton).

% ---------------------------------------------------------------------
% Contributions. ISCA reviewers look for this list. Make each item
% something that can be checked, and point it at its section.
% ---------------------------------------------------------------------
This paper makes the following contributions:
\begin{itemize}
  \item A processor that runs MicroPython~v1.29.0 bytecode as its native ISA, covering
        \textcolor{red}{N} of the 207 defined opcodes with documented
        semantic deviations (Section~\ref{sec:bytecode-isa}).
  \item A type-tagged, ring-buffer register file that maps the
        interpreter's operand stack onto hardware and spills in fixed
        frames (Sections~\ref{sec:regfile}--\ref{sec:tags}).
  \item A collection accelerator and range-based memory and ALU scoreboards
        that overlap container upkeep and arithmetic with execution
        while keeping exceptions precise (Section~\ref{sec:accelerator}).
  \item A hardware/firmware split of the Python runtime, including builtins,
        standard-library modules, garbage collection and a self-hosted
        on-device compiler written in MicroPython (Section~\ref{sec:builtins}).
  \item A post-place-and-route Sky130 implementation evaluated on
        MicroPython's \texttt{perf\_bench} suite against published results
        for five commercial microcontroller configurations (Section~\textcolor{red}{X}).
  % TODO: decide whether the custom benchmarks earn their own bullet.
  % TODO: if the RTL will be open-sourced, say so here (reviewers value
  % artifacts); keep it anonymized for review.
\end{itemize}

% ---------------------------------------------------------------------
% Roadmap. Optional at ISCA; drop it if space is tight.
% ---------------------------------------------------------------------
% Section~\ref{sec:background} reviews how MicroPython executes bytecode
% and derives the requirements for running it natively.
% Section~\ref{sec:design} presents the design of \textmu{}PyCore.
% Sections~X and~Y describe our methodology and results, and
% Section~Z discusses related work.
  (replaces the placeholder
%            \section{Introduction} in the main file)
% Citations: keys resolve against docs/paper/references.bib. To use it,
%            replace the template thebibliography block with
%              \bibliographystyle{IEEEtran}
%              \bibliography{references}
%
% Conventions used below
%   % P<n> GOAL:   what the paragraph must accomplish for a reviewer
%   % EVIDENCE:    which citations carry the claim
%   \textcolor{red}{...}   number still to be finalized (matches the
%                          convention in the Evaluation sections)
%   % TODO:        prose or a fact still to settle
%
% Target length: ~1 to 1.25 columns of text + the contribution list.
% ISCA reviewers read the intro to answer four questions, in order:
%   (1) why does this problem matter now, (2) why do existing
%   solutions fail, (3) what is the key insight, (4) what did you
%   show. Each paragraph below answers one of them.
% =====================================================================

\section{Introduction}
\label{sec:intro}

% ---------------------------------------------------------------------
% P1 GOAL: The "why now". AI coding has moved the bottleneck from
% writing code to reading and reviewing it. Keep it to hard numbers;
% reviewers at ISCA are skeptical of AI hype, so lead with measured
% effects rather than CEO quotes.
% EVIDENCE: pichai2024, nadella2025 (scale of AI-written code);
% faros2025 (review time); becker2025metr (RCT); stackoverflow2025
% (developer distrust -> more review).
% ---------------------------------------------------------------------
Large language models now write a substantial fraction of production
code: Google and Microsoft report that 25\% and 20--30\% of their new
code, respectively, is generated by AI~\cite{pichai2024,nadella2025}.
Writing code is no longer the bottleneck; checking it is. Across
roughly 10,000 developers, teams with high AI adoption merge 98\% more
pull requests, but each one spends 91\% longer in review~\cite{faros2025}.
In a randomized controlled trial, experienced open-source developers
took 19\% \emph{longer} to finish tasks with AI assistance, even though
they believed it had made them faster~\cite{becker2025metr}. Developers
also trust what they review less: 46\% now distrust the accuracy of AI
output, up from 31\% a year earlier~\cite{stackoverflow2025}.
% TODO: one sentence of the authors' thesis: when humans mainly review,
% the language should be optimized for the reader. Python was designed
% around exactly that principle ("Readability counts"~\cite{pep20}).

% ---------------------------------------------------------------------
% P2 GOAL: Python is where human preference and model capability
% converge, so it is the language worth building hardware for. This
% paragraph justifies the choice of Python specifically (a reviewer
% will ask "why not JavaScript / Lua / Java?").
% EVIDENCE: octoverse2024, stackoverflow2025 (humans); twist2025,
% kocetkov2023, cassano2023 (models).
% ---------------------------------------------------------------------
Python is where both trends meet. In 2024 it passed JavaScript as the
most used language on GitHub~\cite{octoverse2024}, and its use among
surveyed developers grew by 7 percentage points in a single
year~\cite{stackoverflow2025}. LLMs show the same preference. Asked to
solve benchmark tasks without being told a language, eight models
chose Python in 90--97\% of cases, and still chose it in 58\% of project
setups where it was a poor fit~\cite{twist2025}. Python is
also heavily represented in the open code corpora these models are
trained on~\cite{kocetkov2023}, and code-generation quality varies
with how well a language is represented in that data~\cite{cassano2023}.
% TODO(verify): before submission, check the direction and size of the
% MultiPL-E popularity/accuracy finding in the paper itself. Search
% snippets confirm only that the paper studies it.

% ---------------------------------------------------------------------
% P3 GOAL: The cost. Python is slow, and the gap is worst exactly
% where a JIT cannot help: memory-constrained embedded devices. This is
% the problem statement; end it with a crisp sentence that a reviewer
% can quote back.
% EVIDENCE: leiserson2020 (47x language gap), hennessy2019 (frames the
% Python gap as an architecture opportunity), ismail2018 (where the
% interpreter's time goes), pereira2021 (energy), micropython,
% eclipse2024 (Python on constrained devices).
% ---------------------------------------------------------------------
That readability comes at a steep cost. Rewriting a Python matrix
multiply in C makes it 47$\times$ faster before any hardware-specific
tuning~\cite{leiserson2020}. In energy-efficiency rankings across 27
languages, Python places near the bottom~\cite{pereira2021}. Hennessy and
Patterson identify this gap as one of the main opportunities for
computer architects~\cite{hennessy2019}. On servers and laptops, JIT
compilers recover part of the gap by spending memory and warm-up
time~\cite{bolz2009pypy,pep744}. Microcontrollers cannot afford either.
MicroPython~\cite{micropython} brings Python to devices with hundreds of
kilobytes of RAM, and 32\% of developers targeting constrained devices
already use Python~\cite{eclipse2024}. Yet these devices run Python
through a bytecode interpreter, where dynamic type dispatch, object
allocation and calls into the C runtime dominate execution
time~\cite{ismail2018}.
% TODO: optionally add one sentence on the size of the MCU market
% (see the yole2025 note in references.bib; a sourced unit count is
% still missing).
% TODO: close P3 with a one-line problem statement, e.g. "Embedded
% developers must choose between the language that people and models
% read best and the performance their device needs."

% ---------------------------------------------------------------------
% P4 GOAL: Why prior work does not solve it. Three buckets, one or
% two sentences each. Each sentence should end with what that bucket
% gives up, which sets up the insight in P5. The detailed versions live
% in Sec. Background / Related Work, so point there.
% EVIDENCE: pep659, pep744, bolz2009pypy (software); micropython
% (native/viper emitters, see Background); mcghan1998, schoeberl2008,
% porthouse2005 (Java processors); moon1985 (tagged Lisp machine);
% kim2017typed, choi2017shortcut, gope2017php, maas2018gc (assists for
% conventional cores); pyxl2025.
% ---------------------------------------------------------------------
Existing approaches recover performance by giving something up.
\emph{Software} techniques such as CPython's specializing
interpreter~\cite{pep659} and tracing or copy-and-patch
JITs~\cite{bolz2009pypy,pep744} need memory a microcontroller lacks.
MicroPython's own native and Viper emitters instead give up the managed
semantics that make the language worth using (Section~\ref{sec:background}).
% TODO: \label{sec:background} does not exist yet; add it to the
% Background section header.
\emph{Language processors} executed bytecode directly, but they targeted
statically typed, verifiable Java bytecode~\cite{mcghan1998,schoeberl2008,porthouse2005}.
They kept the interpreter's hard cases, such as object operations, in
software. Earlier tagged Lisp machines~\cite{moon1985} showed that
type tags in hardware help a dynamic language. \emph{Architectural
assists} for scripting languages~\cite{kim2017typed,choi2017shortcut,gope2017php,maas2018gc}
add type checks, dispatch shortcuts, or hash and GC
engines to a conventional core that still runs an interpreter or JIT.
The reported gains are 10--30\%.
% TODO(verify): the 10-30% range combines kim2017typed (about 10-11%
% geomean) and choi2017shortcut (30% under the baseline compiler).
% Check both against the papers, or drop the number.
Most recently, PyXL~\cite{pyxl2025} runs Python on an FPGA, but it
compiles to a custom instruction set on a host machine and is not
publicly documented.

% ---------------------------------------------------------------------
% P5 GOAL: The key insight + what µPyCore is. State the thesis in one
% sentence, then map each obstacle (from the Background "rules") to the
% mechanism that removes it. This is the paragraph reviewers summarize
% in their review, so make the mapping explicit.
% ---------------------------------------------------------------------
This paper takes the opposite approach: rather than lowering Python
toward the machine, we raise the machine to Python. \textmu{}PyCore
executes unmodified MicroPython bytecode as its native instruction set
and implements in hardware the runtime machinery that a conventional
core spends most of its time emulating. Three obstacles shape the
design.
% TODO: tighten each to one sentence once Sec. Design is final.
First, bytecode carries no type information, so every value in
\textmu{}PyCore carries a 4-bit type tag that the core reads together
with its payload in a single access (Section~\ref{sec:tags}).
Second, the bytecode targets a stack machine, so the register file is a
tagged ring buffer that holds the top of the execution stack and
spills to SRAM in fixed frames (Section~\ref{sec:regfile}).
Third, the language's time goes to built-in containers, allocation and
garbage collection rather than to arithmetic. A collection accelerator
and two range-based scoreboards therefore run that work in parallel with
the core (Sections~\ref{sec:accelerator}--\ref{sec:builtins}).
% TODO: Scoreboarding subsection needs a \label (e.g. sec:scoreboard).
Everything the hardware does not cover, including \texttt{print},
\texttt{sorted} and a full bytecode compiler, runs as MicroPython
firmware on the same core, so Python source can be compiled and run on
the device itself.

% ---------------------------------------------------------------------
% P6 GOAL: Headline results. Same numbers as the abstract; keep red
% until the evaluation is frozen. State the methodology caveat
% (RTL simulation at the post-P&R clock vs. published board numbers)
% in the same breath so a reviewer cannot say it was hidden.
% ---------------------------------------------------------------------
We synthesize and place-and-route \textmu{}PyCore in SkyWater's open
130-nm process~\cite{skywater130}, where it closes timing at
\textcolor{red}{83.33~MHz}. Cycle counts from RTL simulation at that
clock project \textcolor{red}{7,286}~pystones/s,
\textcolor{red}{2.84$\times$} the published throughput of an Arm-based
Raspberry Pi Pico~2 at 150~MHz and \textcolor{red}{3.19$\times$} that of a
PyBoard at 168~MHz. Across \textcolor{red}{15} workloads from
MicroPython's \texttt{perf\_bench} suite, \textmu{}PyCore achieves a
geometric-mean speedup of \textcolor{red}{2.33$\times$} over the fastest
published microcontroller result for each workload, while running at a
lower clock rate than any of them.
% TODO: reconcile the numbers with the abstract and Sec. Eval before
% submission (see the notes left with this skeleton).

% ---------------------------------------------------------------------
% Contributions. ISCA reviewers look for this list. Make each item
% something that can be checked, and point it at its section.
% ---------------------------------------------------------------------
This paper makes the following contributions:
\begin{itemize}
  \item A processor that runs MicroPython~v1.29.0 bytecode as its native ISA, covering
        \textcolor{red}{N} of the 207 defined opcodes with documented
        semantic deviations (Section~\ref{sec:bytecode-isa}).
  \item A type-tagged, ring-buffer register file that maps the
        interpreter's operand stack onto hardware and spills in fixed
        frames (Sections~\ref{sec:regfile}--\ref{sec:tags}).
  \item A collection accelerator and range-based memory and ALU scoreboards
        that overlap container upkeep and arithmetic with execution
        while keeping exceptions precise (Section~\ref{sec:accelerator}).
  \item A hardware/firmware split of the Python runtime, including builtins,
        standard-library modules, garbage collection and a self-hosted
        on-device compiler written in MicroPython (Section~\ref{sec:builtins}).
  \item A post-place-and-route Sky130 implementation evaluated on
        MicroPython's \texttt{perf\_bench} suite against published results
        for five commercial microcontroller configurations (Section~\textcolor{red}{X}).
  % TODO: decide whether the custom benchmarks earn their own bullet.
  % TODO: if the RTL will be open-sourced, say so here (reviewers value
  % artifacts); keep it anonymized for review.
\end{itemize}

% ---------------------------------------------------------------------
% Roadmap. Optional at ISCA; drop it if space is tight.
% ---------------------------------------------------------------------
% Section~\ref{sec:background} reviews how MicroPython executes bytecode
% and derives the requirements for running it natively.
% Section~\ref{sec:design} presents the design of \textmu{}PyCore.
% Sections~X and~Y describe our methodology and results, and
% Section~Z discusses related work.
  (replaces the placeholder
%            \section{Introduction} in the main file)
% Citations: keys resolve against docs/paper/references.bib. To use it,
%            replace the template thebibliography block with
%              \bibliographystyle{IEEEtran}
%              \bibliography{references}
%
% Conventions used below
%   % P<n> GOAL:   what the paragraph must accomplish for a reviewer
%   % EVIDENCE:    which citations carry the claim
%   \textcolor{red}{...}   number still to be finalized (matches the
%                          convention in the Evaluation sections)
%   % TODO:        prose or a fact still to settle
%
% Target length: ~1 to 1.25 columns of text + the contribution list.
% ISCA reviewers read the intro to answer four questions, in order:
%   (1) why does this problem matter now, (2) why do existing
%   solutions fail, (3) what is the key insight, (4) what did you
%   show. Each paragraph below answers one of them.
% =====================================================================

\section{Introduction}
\label{sec:intro}

% ---------------------------------------------------------------------
% P1 GOAL: The "why now". AI coding has moved the bottleneck from
% writing code to reading and reviewing it. Keep it to hard numbers;
% reviewers at ISCA are skeptical of AI hype, so lead with measured
% effects rather than CEO quotes.
% EVIDENCE: pichai2024, nadella2025 (scale of AI-written code);
% faros2025 (review time); becker2025metr (RCT); stackoverflow2025
% (developer distrust -> more review).
% ---------------------------------------------------------------------
Large language models now write a substantial fraction of production
code: Google and Microsoft report that 25\% and 20--30\% of their new
code, respectively, is generated by AI~\cite{pichai2024,nadella2025}.
Writing code is no longer the bottleneck; checking it is. Across
roughly 10,000 developers, teams with high AI adoption merge 98\% more
pull requests, but each one spends 91\% longer in review~\cite{faros2025}.
In a randomized controlled trial, experienced open-source developers
took 19\% \emph{longer} to finish tasks with AI assistance, even though
they believed it had made them faster~\cite{becker2025metr}. Developers
also trust what they review less: 46\% now distrust the accuracy of AI
output, up from 31\% a year earlier~\cite{stackoverflow2025}.
% TODO: one sentence of the authors' thesis: when humans mainly review,
% the language should be optimized for the reader. Python was designed
% around exactly that principle ("Readability counts"~\cite{pep20}).

% ---------------------------------------------------------------------
% P2 GOAL: Python is where human preference and model capability
% converge, so it is the language worth building hardware for. This
% paragraph justifies the choice of Python specifically (a reviewer
% will ask "why not JavaScript / Lua / Java?").
% EVIDENCE: octoverse2024, stackoverflow2025 (humans); twist2025,
% kocetkov2023, cassano2023 (models).
% ---------------------------------------------------------------------
Python is where both trends meet. In 2024 it passed JavaScript as the
most used language on GitHub~\cite{octoverse2024}, and its use among
surveyed developers grew by 7 percentage points in a single
year~\cite{stackoverflow2025}. LLMs show the same preference. Asked to
solve benchmark tasks without being told a language, eight models
chose Python in 90--97\% of cases, and still chose it in 58\% of project
setups where it was a poor fit~\cite{twist2025}. Python is
also heavily represented in the open code corpora these models are
trained on~\cite{kocetkov2023}, and code-generation quality varies
with how well a language is represented in that data~\cite{cassano2023}.
% TODO(verify): before submission, check the direction and size of the
% MultiPL-E popularity/accuracy finding in the paper itself. Search
% snippets confirm only that the paper studies it.

% ---------------------------------------------------------------------
% P3 GOAL: The cost. Python is slow, and the gap is worst exactly
% where a JIT cannot help: memory-constrained embedded devices. This is
% the problem statement; end it with a crisp sentence that a reviewer
% can quote back.
% EVIDENCE: leiserson2020 (47x language gap), hennessy2019 (frames the
% Python gap as an architecture opportunity), ismail2018 (where the
% interpreter's time goes), pereira2021 (energy), micropython,
% eclipse2024 (Python on constrained devices).
% ---------------------------------------------------------------------
That readability comes at a steep cost. Rewriting a Python matrix
multiply in C makes it 47$\times$ faster before any hardware-specific
tuning~\cite{leiserson2020}. In energy-efficiency rankings across 27
languages, Python places near the bottom~\cite{pereira2021}. Hennessy and
Patterson identify this gap as one of the main opportunities for
computer architects~\cite{hennessy2019}. On servers and laptops, JIT
compilers recover part of the gap by spending memory and warm-up
time~\cite{bolz2009pypy,pep744}. Microcontrollers cannot afford either.
MicroPython~\cite{micropython} brings Python to devices with hundreds of
kilobytes of RAM, and 32\% of developers targeting constrained devices
already use Python~\cite{eclipse2024}. Yet these devices run Python
through a bytecode interpreter, where dynamic type dispatch, object
allocation and calls into the C runtime dominate execution
time~\cite{ismail2018}.
% TODO: optionally add one sentence on the size of the MCU market
% (see the yole2025 note in references.bib; a sourced unit count is
% still missing).
% TODO: close P3 with a one-line problem statement, e.g. "Embedded
% developers must choose between the language that people and models
% read best and the performance their device needs."

% ---------------------------------------------------------------------
% P4 GOAL: Why prior work does not solve it. Three buckets, one or
% two sentences each. Each sentence should end with what that bucket
% gives up, which sets up the insight in P5. The detailed versions live
% in Sec. Background / Related Work, so point there.
% EVIDENCE: pep659, pep744, bolz2009pypy (software); micropython
% (native/viper emitters, see Background); mcghan1998, schoeberl2008,
% porthouse2005 (Java processors); moon1985 (tagged Lisp machine);
% kim2017typed, choi2017shortcut, gope2017php, maas2018gc (assists for
% conventional cores); pyxl2025.
% ---------------------------------------------------------------------
Existing approaches recover performance by giving something up.
\emph{Software} techniques such as CPython's specializing
interpreter~\cite{pep659} and tracing or copy-and-patch
JITs~\cite{bolz2009pypy,pep744} need memory a microcontroller lacks.
MicroPython's own native and Viper emitters instead give up the managed
semantics that make the language worth using (Section~\ref{sec:background}).
% TODO: \label{sec:background} does not exist yet; add it to the
% Background section header.
\emph{Language processors} executed bytecode directly, but they targeted
statically typed, verifiable Java bytecode~\cite{mcghan1998,schoeberl2008,porthouse2005}.
They kept the interpreter's hard cases, such as object operations, in
software. Earlier tagged Lisp machines~\cite{moon1985} showed that
type tags in hardware help a dynamic language. \emph{Architectural
assists} for scripting languages~\cite{kim2017typed,choi2017shortcut,gope2017php,maas2018gc}
add type checks, dispatch shortcuts, or hash and GC
engines to a conventional core that still runs an interpreter or JIT.
The reported gains are 10--30\%.
% TODO(verify): the 10-30% range combines kim2017typed (about 10-11%
% geomean) and choi2017shortcut (30% under the baseline compiler).
% Check both against the papers, or drop the number.
Most recently, PyXL~\cite{pyxl2025} runs Python on an FPGA, but it
compiles to a custom instruction set on a host machine and is not
publicly documented.

% ---------------------------------------------------------------------
% P5 GOAL: The key insight + what µPyCore is. State the thesis in one
% sentence, then map each obstacle (from the Background "rules") to the
% mechanism that removes it. This is the paragraph reviewers summarize
% in their review, so make the mapping explicit.
% ---------------------------------------------------------------------
This paper takes the opposite approach: rather than lowering Python
toward the machine, we raise the machine to Python. \textmu{}PyCore
executes unmodified MicroPython bytecode as its native instruction set
and implements in hardware the runtime machinery that a conventional
core spends most of its time emulating. Three obstacles shape the
design.
% TODO: tighten each to one sentence once Sec. Design is final.
First, bytecode carries no type information, so every value in
\textmu{}PyCore carries a 4-bit type tag that the core reads together
with its payload in a single access (Section~\ref{sec:tags}).
Second, the bytecode targets a stack machine, so the register file is a
tagged ring buffer that holds the top of the execution stack and
spills to SRAM in fixed frames (Section~\ref{sec:regfile}).
Third, the language's time goes to built-in containers, allocation and
garbage collection rather than to arithmetic. A collection accelerator
and two range-based scoreboards therefore run that work in parallel with
the core (Sections~\ref{sec:accelerator}--\ref{sec:builtins}).
% TODO: Scoreboarding subsection needs a \label (e.g. sec:scoreboard).
Everything the hardware does not cover, including \texttt{print},
\texttt{sorted} and a full bytecode compiler, runs as MicroPython
firmware on the same core, so Python source can be compiled and run on
the device itself.

% ---------------------------------------------------------------------
% P6 GOAL: Headline results. Same numbers as the abstract; keep red
% until the evaluation is frozen. State the methodology caveat
% (RTL simulation at the post-P&R clock vs. published board numbers)
% in the same breath so a reviewer cannot say it was hidden.
% ---------------------------------------------------------------------
We synthesize and place-and-route \textmu{}PyCore in SkyWater's open
130-nm process~\cite{skywater130}, where it closes timing at
\textcolor{red}{83.33~MHz}. Cycle counts from RTL simulation at that
clock project \textcolor{red}{7,286}~pystones/s,
\textcolor{red}{2.84$\times$} the published throughput of an Arm-based
Raspberry Pi Pico~2 at 150~MHz and \textcolor{red}{3.19$\times$} that of a
PyBoard at 168~MHz. Across \textcolor{red}{15} workloads from
MicroPython's \texttt{perf\_bench} suite, \textmu{}PyCore achieves a
geometric-mean speedup of \textcolor{red}{2.33$\times$} over the fastest
published microcontroller result for each workload, while running at a
lower clock rate than any of them.
% TODO: reconcile the numbers with the abstract and Sec. Eval before
% submission (see the notes left with this skeleton).

% ---------------------------------------------------------------------
% Contributions. ISCA reviewers look for this list. Make each item
% something that can be checked, and point it at its section.
% ---------------------------------------------------------------------
This paper makes the following contributions:
\begin{itemize}
  \item A processor that runs MicroPython~v1.29.0 bytecode as its native ISA, covering
        \textcolor{red}{N} of the 207 defined opcodes with documented
        semantic deviations (Section~\ref{sec:bytecode-isa}).
  \item A type-tagged, ring-buffer register file that maps the
        interpreter's operand stack onto hardware and spills in fixed
        frames (Sections~\ref{sec:regfile}--\ref{sec:tags}).
  \item A collection accelerator and range-based memory and ALU scoreboards
        that overlap container upkeep and arithmetic with execution
        while keeping exceptions precise (Section~\ref{sec:accelerator}).
  \item A hardware/firmware split of the Python runtime, including builtins,
        standard-library modules, garbage collection and a self-hosted
        on-device compiler written in MicroPython (Section~\ref{sec:builtins}).
  \item A post-place-and-route Sky130 implementation evaluated on
        MicroPython's \texttt{perf\_bench} suite against published results
        for five commercial microcontroller configurations (Section~\textcolor{red}{X}).
  % TODO: decide whether the custom benchmarks earn their own bullet.
  % TODO: if the RTL will be open-sourced, say so here (reviewers value
  % artifacts); keep it anonymized for review.
\end{itemize}

% ---------------------------------------------------------------------
% Roadmap. Optional at ISCA; drop it if space is tight.
% ---------------------------------------------------------------------
% Section~\ref{sec:background} reviews how MicroPython executes bytecode
% and derives the requirements for running it natively.
% Section~\ref{sec:design} presents the design of \textmu{}PyCore.
% Sections~X and~Y describe our methodology and results, and
% Section~Z discusses related work.
